% Vietnamese translation for OI Public Library
% Source-File: IOI中国国家候选队论文/2007/day2/6.何森《浅谈数据的合理组织》.ppt
\documentclass[11pt]{article}
\usepackage{oipl-vn}

\newcommand{\Oh}{\mathcal{O}}

\title{Bản trình chiếu: Cách tổ chức dữ liệu hợp lý}
\author{He Sen}
\date{2007}

\begin{document}
\maketitle

\origtitle{Slide companion for reasonable organization of data}
\sourcepath{IOI China candidate team papers / 2007 / day 2 / He Sen.ppt}

\section{Phạm vi bản dịch}

Tệp PowerPoint đi cùng bài viết đã được dịch từ PDF. Bản ghi chú này giữ lại
hai góc nhìn chính: tổ chức lại cấu trúc dữ liệu và tổ chức lại thứ tự dữ
liệu để làm rõ mô hình hoặc tăng hiệu quả thuật toán.

\section{Hai hướng tổ chức}

Với cùng một tập dữ liệu, cách lưu trữ khác nhau có thể tạo ra thuật toán rất
khác nhau. Bảng băm và trie thay đổi hình thức biểu diễn; sắp xếp trước rồi
tìm kiếm nhị phân thay đổi thứ tự xử lý. Trong bài phức tạp, hai hướng này
thường được phối hợp.

\section{Ví dụ quy hoạch động}

Trong bài mua sắm có món chính và phụ kiện, nếu nhìn dữ liệu như một cây tổng
quát, ta có thể viết DP trên cây nhưng độ phức tạp cao. Quan sát thêm rằng
mỗi món chính có nhiều nhất hai phụ kiện và phụ kiện không có con, ta gom mỗi
món chính cùng phụ kiện thành một nhóm lựa chọn nhỏ. Khi đó bài toán trở lại
gần với ba lô \(0\)-\(1\), đơn giản và nhanh hơn.

\section{Vai trò trong mô hình hóa}

Tổ chức dữ liệu không chỉ là chi tiết cài đặt. Khi đặt lại quan hệ cha con,
thứ tự duyệt, nhóm đối tượng hoặc lớp tương đương, ta có thể phát hiện chuyển
trạng thái ngắn hơn, giảm số chiều DP, hoặc biến đồ thị khó xử lý thành một
cấu trúc quen thuộc như cây, dãy hoặc các nhóm độc lập.

\end{document}
